% paper/oe_main.tex -- Phase 7.2 manuscript skeleton for Optics Express.
%
% ============================================================================
% STATUS: SKELETON, NOT SUBMITTABLE. Read this block before editing further.
% ============================================================================
% 1. TEMPLATE: this uses a generic article-like structure, NOT Optica's
%    actual class file. Per the master prompt's own instruction ("the user
%    downloads the template and confirms the current class file and
%    abstract word limit from the author guidelines -- do not guess
%    these"), I have not fetched or guessed Optica's template (no web
%    access in this environment). Once you've downloaded the real
%    opticajnl.cls (or current equivalent) and confirmed the abstract word
%    limit, swap \documentclass and re-flow content into its required
%    structure.
% 2. NUMBERS: every quantitative claim below is a \macro from numbers.tex
%    (scripts/make_numbers_tex.py, generated from
%    results/summary/paper_numbers.json). Never hand-type a number here --
%    scripts/check_consistency.py fails the build if you do. As of this
%    pass, Phase 3 (E1-E6 GPU runs) has not happened, so most macros
%    render as [PENDING] -- this file is correct-but-incomplete by
%    construction until real data lands, not broken.
% 3. arXiv:2601.16330 VERIFIED (via WebSearch/WebFetch): resolves to
%    Wechsler, Rizzo, Moser, "Single-View Holographic Volumetric 3D
%    Printing with Coupled Differentiable Wave-Optical and Photochemical
%    Optimization" (SHVAM), submitted 2026-01-22. Confirmed: couples a
%    differentiable wave-optical forward model with a photochemical
%    model, as the draft's differentiators paragraph assumed. NUANCE TO
%    CHECK BEFORE ASSERTING THE DIFFERENTIATORS AS WRITTEN: its
%    photochemical model "explicitly captures inhibitor diffusion" per
%    its own abstract -- i.e. it does include a transport/diffusion term,
%    contrary to a differentiator claim of "no diffusion/transport
%    term at all." The correct differentiator is almost certainly still
%    real (printed/threshold-converted geometry vs. a recorded Bragg-
%    selective index grating; likely no NON-LOCAL polymerization kernel
%    specifically, which is different from having no diffusion term at
%    all) but was NOT independently confirmed here -- read the actual
%    PDF before finalizing this paragraph's wording, don't assert the
%    stronger "no diffusion" version.
% 4. check_refs.py: 3 of the original 4 TODO-VERIFY bibliography entries
%    are now resolved with verified titles/authors/venues (gleeson2008chain
%    confirmed correct as originally cited; kelly2011reply and
%    jeong2022pqpmma were not just missing titles but WRONG -- wrong
%    author, year, and/or volume/page -- corrected and renamed to
%    sheridan2012reply / hsieh2022pqpmma respectively, see paper/refs.bib
%    for the full trail). gallego2008shrinkage remains unresolved: its
%    cited PubMed ID resolves to a different paper entirely (different
%    author, journal, and year) -- see the comment in refs.bib. This file
%    that's resolved, by design.
% ============================================================================

\documentclass[10pt]{article}
\usepackage[margin=1in]{geometry}
\usepackage{amsmath,amssymb,graphicx,booktabs,hyperref,cite}
% AUTO-GENERATED by scripts/make_numbers_tex.py -- DO NOT HAND-EDIT.
% Regenerate: python scripts/make_numbers_tex.py
\usepackage{xcolor}  % for the [PENDING] flag color, harmless if already loaded

\newcommand{\MeasuredContrastTwoX}{2.03}
\newcommand{\KcPredTwoX}{4.25}
\newcommand{\KstarTwoXInterp}{\textcolor{red}{[PENDING]}}
\newcommand{\KstarTwoXCI}{7.85}
\newcommand{\MeanGainTwoX}{0.34}
\newcommand{\MeasuredContrastFourX}{3.98}
\newcommand{\KcPredFourX}{6.96}
\newcommand{\KstarFourXInterp}{\textcolor{red}{[PENDING]}}
\newcommand{\KstarFourXCI}{15.71}
\newcommand{\MeanGainFourX}{0.65}
\newcommand{\MeasuredContrastEightX}{6.03}
\newcommand{\KcPredEightX}{7.39}
\newcommand{\KstarEightXInterp}{\textcolor{red}{[PENDING]}}
\newcommand{\KstarEightXCI}{15.71}
\newcommand{\MeanGainEightX}{0.73}
\newcommand{\SeedCount}{\textcolor{red}{[PENDING]}}
\newcommand{\NResultFiles}{684}

% --- supplement macros (already-real data, not Phase-3-gated) ---
\newcommand{\AblationCheckpointCosSim}{0.977}
\newcommand{\AblationSurrogateCosSim}{0.841}
\newcommand{\RCWAMaxDevThreeCase}{0.038}
\newcommand{\RCWAE7MaxDev}{0.980}
\newcommand{\RCWAE7NCases}{90}
\newcommand{\MeshPSNRSpread}{0.036}
\newcommand{\WavelengthDesignPSNR}{6.63}
  % scripts/make_numbers_tex.py output -- all quantitative macros

\title{Media-in-the-Loop Holography: Computer-Generated Holography\\
Through a Differentiable Model of Volume Photopolymer Recording}
\author{Chaitanya Tripathi$^1$, [Co-author]$^2$\\
\small $^1$Ashoka University, Sonipat, India \quad $^2$BITS Pilani, India}
\date{}

\begin{document}
\maketitle

\begin{abstract}
% TODO (blocked on Phase 3): master prompt spec is "<=~100 words (pending
% guideline check), leading with the budget-controlled cliff finding."
% The budget-controlled cliff finding does not exist yet -- E1 (the
% experiment that tests whether the cliff location tracks budget, per
% Eq. 5) has not run. Writing a lead sentence asserting an outcome before
% the experiment exists would be exactly the "future work described as
% done" the master prompt's Definition of Done forbids. Structural
% skeleton below uses macros; fill in the actual finding once E1's
% headroom-closure table (analysis/aggregate.py) has real numbers.
Computer-generated holography algorithms overwhelmingly assume an ideal
phase modulator; camera-in-the-loop correction is structurally unavailable
for write-once volume photopolymer recording media. We introduce
media-in-the-loop holography: a differentiable digital twin of the NPDD
recording chemistry placed inside the hologram-design loop. Across
\SeedCount\ seeds and three contrast-headroom budgets, media-in-the-loop
optimization recovers \MeanGainTwoX{}/\MeanGainFourX{}/\MeanGainEightX{}~dB
over media-blind SGD at $2\times/4\times/8\times$ budget respectively
[PENDING E1], and the observed compensation-cliff location
$K^*=\KstarTwoXInterp$ [PENDING] tracks the analytic prediction
$K_c(\text{measured } C)=\KcPredTwoX$ [PENDING] as the budget moves. All
code, manifests, and result archives are released.
\end{abstract}

\section{Introduction}
Phase-only CGH computes a modulation pattern whose diffracted field
approximates a target image. Camera-in-the-loop (CITL) optimization
closed the gap between the computed pattern and physical spatial light
modulator hardware by measuring the actual reconstruction and
backpropagating the residual into the display model~\cite{peng2020neural}.

This solution is structurally unavailable to an entire class of
holographic devices. Volume photopolymer media -- Bayfol HX,
PVA/acrylamide, PQ/PMMA -- record a hologram as a permanent
refractive-index distribution formed by photopolymerization and mass
transport. The medium is write-once: measurement of the reconstruction
error is possible only after the error has become unerasable. CITL's
feedback signal arrives, by construction, too late. Yet these media are
precisely where model mismatch is most severe: the recorded index
profile is not a scaled copy of the exposure, but the product of monomer
diffusion, non-local polymer chain growth, dye depletion, saturable index
response, and post-exposure shrinkage, each depending nonlinearly on the
exposure itself.

We close this gap by placing a differentiable model of the recording
chemistry -- not just the readout optics -- inside the design loop.

\section{Related Work}

\textbf{Model-mismatch correction for holographic displays.} CITL
optimization~\cite{peng2020neural} and its extensions correct the
display's forward model using camera measurements; all require a
rewritable modulator observed live. Our setting inverts the
availability: the medium is write-once, so the model must carry the
full burden.

\textbf{Fabrication-aware surface-relief optics.} Neural
lithography~\cite{zheng2023neural} inserts a learned digital twin of a
lithographic process into differentiable design, but models surface
geometry -- an etched height map. Volume photopolymers differ in kind:
the recorded object is a volumetric index distribution produced by
chemistry, and readout is Bragg-selective.

% VERIFIED (see file header note 3): arXiv:2601.16330 is Wechsler, Rizzo,
% Moser, "Single-View Holographic Volumetric 3D Printing with Coupled
% Differentiable Wave-Optical and Photochemical Optimization" (SHVAM),
% 2026. Differentiators per the master prompt, checked against the
% verified abstract: printed/threshold-converted geometry vs. our
% recorded, Bragg-selective index grating -- holds. "No compensation-
% limit theory" -- not contradicted by the abstract, but not positively
% confirmed either; SHVAM's own framing is dose-distribution shaping, not
% a K-dependent compensation-cliff argument. "No NPDD non-locality/
% transport term" needs REWORDING, not deletion: SHVAM's abstract states
% its photochemical model "explicitly captures inhibitor diffusion," so
% it does have a diffusion/transport term -- the accurate differentiator
% is no NON-LOCAL (spatially-smeared production) kernel specifically,
% which is a different and narrower claim. Read the full PDF before
% finalizing this paragraph's exact wording.
\textbf{[TODO: write the concurrent-work paragraph -- citation and
differentiators are now verified (see comment above), narrow prose
still needs drafting, particularly the diffusion/non-locality distinction.]}
~\cite{wechsler2026shvam}

% TODO: (b) adjoint/PDE-constrained inverse design in photonics -- add
% citations once selected (not fabricated here).
% TODO: (c) proximity-effect correction in e-beam lithography as the
% conceptual cousin of the compensation cliff -- add citation.
% TODO: (d) NPDD primary sources Sheridan & Lawrence 2000, Zhao &
% Mouroulis 1994 -- exact records to be verified by the user (per the
% master prompt's own instruction), not fabricated here.

Our novelty claim is scoped precisely: the first differentiable
\emph{recording} model for volume holographic media placed inside a
hologram-design loop -- not the first differentiable photochemistry
model in any loop.

\section{Model}
The twin couples the non-local polymerization-driven diffusion (NPDD)
reaction-diffusion system to volume diffraction via split-step beam
propagation, with gradients obtained by unrolling a spectral IMEX
integrator. One modeling note: the non-local term is
$F\cdot(G_\sigma * u)$, not $G_\sigma * (F\cdot u)$ -- the kernel must
smear the full production term, since chains initiated at $x'$ deposit
polymer at $x$ (Sheridan NPDD formulation).

\section{Media-in-the-Loop Optimization}
Given target intensity $I_t$, we optimize the nonnegative, dose-budgeted
exposure $E(x)$ delivered to the medium, with the twin inside the loss.
Contrast headroom is enforced as a hard constraint
(\texttt{holomedia.optimize.contrast\_project}), not merely a nominal
dose level -- see docs/definitions.md (d) for the audit that found no
such constraint existed prior to this pass.

\section{Results}
% TODO (blocked on Phase 3): "multi-seed numbers only" per the master
% prompt. Cliff table, method comparison (M1-M5b), and the K*-vs-Kc
% headroom-closure table all require E1 to have actually run. Structural
% placeholders below.
\textbf{[TODO: cliff table sorted by K, GS and linear-precomp columns, mean $\pm$ std -- pending E1.]}

\textbf{[TODO: Figs. F4 (headline gain vs K) and F5 (K* vs Kc scatter) -- pending E1; see figures/make\_all.py, currently emitting placeholders for these.]}

\section{Twin Validation}
% TODO (blocked on Phase 6): pending digitized literature CSVs and
% experiments/fit_literature_curves.py's fit-quality report. Report the
% ACTUAL fit quality (good or bad) once real data exists -- do not assert
% agreement here.
\textbf{[TODO: twin-vs-literature overlay and fit quality -- pending Phase 6 (data/literature/*.csv).]}

\section{Discussion and Limitations}
This study is simulation-only; its claims are conditioned on a physics
model (NPDD, Kogelnik, split-step BPM) validated component-wise
(Sec.~\ref{sec:twin-validation-ref}) but not yet against a completed
literature overlay. Scalar diffraction bounds validity to moderate
numerical apertures; the RCWA validity-envelope grid (E7) bounds this
quantitatively for the K/$\Delta n$/polarization/slant range tested --
readout-side error moves absolute PSNRs more than paired gains, since the
compensation cliff is a recording-side phenomenon (Eq.~5). Recording
kinetics are 1D in transverse $x$; depth ($z$) enters only at readout via
the split-step propagation, not as depth-resolved chemistry -- see
docs/definitions.md (c) for the precise statement of this geometry, which
an earlier draft left ambiguous. Experimental two-beam closure (recording
optimized vs. naive exposures in a characterized medium and comparing
measured reconstructions) is explicitly future work.

\section{Code and Data Availability}
All source code, experiment manifests, and raw result JSONs are available
at \url{https://github.com/ultramagnus23/HoloForge} under
\texttt{holography\_media/}, released under the Apache 2.0 license.
% TODO (Phase 7.5, blocked on results freeze): a versioned Zenodo DOI
% archive is prepared at results freeze -- not yet cut, since Phase 3
% results don't exist yet to freeze. See REPRODUCE.md for the exact
% regeneration commands once results land.
\textbf{[TODO: Zenodo DOI -- pending results freeze, see REPRODUCE.md.]}

\bibliographystyle{plain}
\bibliography{refs}

\end{document}
